\documentclass[11pt]{article}
\usepackage[margin=1in]{geometry}
\usepackage{amsmath,amssymb,amsthm}
\usepackage{graphicx}
\usepackage[hidelinks]{hyperref}

\newtheorem{theorem}{Theorem}
\newtheorem{lemma}{Lemma}
\newtheorem{remark}{Remark}

\newcommand{\D}{\mathbb D}
\newcommand{\Aut}{\operatorname{Aut}(\D)}
\newcommand{\HH}{\mathcal H}
\newcommand{\BMOA}{\operatorname{BMOA}}
\newcommand{\MD}{M(\mathcal D_\mu)}
\newcommand{\dA}{\,dA}

\title{The non-radial Hilbert-matrix norm problem}
\author{arXiv automatic math run \texttt{fa\_banach\_001}}
\date{2026-06-18}

\begin{document}
\maketitle

\section*{Source and question}

The source paper is Bellavita--Stylogiannis, \emph{Hilbert matrix operator
Acting between conformally invariant spaces}, arXiv:2408.01060
\cite{BellavitaStylogiannis}.  Its final section records the following
question.

\begin{center}
\includegraphics[width=0.95\linewidth]{figures/open_problem_crop.png}
\end{center}

For a positive Borel measure \(\mu\) on \(\D\), the source assumes
\[
 \int_{\D}(1-|w|^2)\,d\mu(w)<\infty
\]
and sets
\[
 U_\mu(z)=\int_{\D}\log\left|\frac{1-\overline w z}{z-w}\right|^2\,d\mu(w).
\]
The norm in \(M(\mathcal D_\mu)\) is
\[
 \|g\|_{\MD}
 =
 |g(0)|+
 \sup_{\phi\in\Aut}
 \left(\int_{\D}|g'(z)|^2 U_\mu(\phi(z))\dA(z)\right)^{1/2}.
\]

\section*{Result}

\begin{theorem}
Let \(\mu\) be a nonzero positive Borel measure on \(\D\) satisfying
\[
 \int_{\D}(1-|w|^2)\,d\mu(w)<\infty .
\]
Then \(\HH:H^\infty\to\MD\) is bounded if and only if
\(\mu(\D)<\infty\).  In the bounded case,
\[
 \|\HH\|_{H^\infty\to \MD}
 =
 1+\frac{\pi}{\sqrt2}\sqrt{\mu(\D)}.
\]
Consequently Question 8.1 has the above exact answer for arbitrary non-radial
\(\mu\) in the source class: finite mass gives the displayed norm, while
infinite mass gives unboundedness.
\end{theorem}

\section*{Proof intuition}

The \(M(\mathcal D_\mu)\)-seminorm is an average, over \(w\), of usual local
\(\BMOA\) Green seminorms after applying an automorphism.  A degenerating
automorphism \(\sigma_r\), \(r\uparrow1\), sends every fixed point \(w\in\D\)
to the same boundary point.  For both \(L(z)=\log(1/(1-z))\) and
\(K(z)=\HH(1)(z)\), the local \(\BMOA\) seminorm at \(\sigma_r(w)\) tends to
the extremal value \(\pi/\sqrt2\).

If \(\mu\) has infinite total mass, apply this boundary degeneration on
finite-measure truncations of arbitrarily large mass.  The seminorm of
\(\HH(1)\) is forced to be infinite, so \(\HH\) is not bounded.  If \(\mu\) is
finite, dominated convergence gives the exact seminorm
\((\pi/\sqrt2)\sqrt{\mu(\D)}\), and the standard derivative identity
\((1-z)(\HH f)'=b_f\), \(\|b_f\|_\infty\leq\|f\|_\infty\), gives the sharp upper
bound.  Testing \(f\equiv1\) gives equality.

\section*{Proof}

Write
\[
 G(z,w)=\log\left|\frac{1-\overline w z}{z-w}\right|^2
\]
and define the local Green seminorm
\[
 \mathcal D_a(g)=\int_{\D}|g'(z)|^2G(z,a)\dA(z),\qquad a\in\D .
\]
Green's function is invariant under automorphisms:
\[
 G(\phi(z),w)=G(z,\phi^{-1}(w)),\qquad \phi\in\Aut .
\]
Therefore Tonelli's theorem gives, for every analytic \(g\) for which either
side is finite,
\begin{equation}\label{eq:green-average}
\int_{\D}|g'(z)|^2U_\mu(\phi(z))\dA(z)
=
\int_{\D}\mathcal D_{\phi^{-1}(w)}(g)\,d\mu(w).
\end{equation}

We shall use the two boundary-local facts already implicit in the source
paper's proof of its \(H^\infty\to\BMOA\) norm estimate.  Let
\[
 L(z)=\log\frac1{1-z},\qquad
 K(z)=\HH(1)(z)=\frac1z\log\frac1{1-z}
\]
with \(K(0)=1\).  Then
\begin{equation}\label{eq:local-limits}
 \sup_{a\in\D}\mathcal D_a(L)^{1/2}
 =
 \sup_{a\in\D}\mathcal D_a(K)^{1/2}
 =
 \frac{\pi}{\sqrt2},
\end{equation}
and, for every fixed \(w\in\D\),
\begin{equation}\label{eq:boundary-limits}
 \mathcal D_{\sigma_r(w)}(L)^{1/2}\to\frac{\pi}{\sqrt2},
 \qquad
 \mathcal D_{\sigma_r(w)}(K)^{1/2}\to\frac{\pi}{\sqrt2}
\end{equation}
as \(r\uparrow1\), where \(\sigma_r(z)=(r-z)/(1-rz)\).  For \(K\), this is the
source's decomposition
\[
 zK'(z)=-\frac1z\log\frac1{1-z}+\frac1{1-z};
\]
the logarithmic term is negligible in the degenerating local Green seminorm,
and the second term has the same boundary limit as \(L'\).

First suppose that \(\mu(\D)=\infty\).  The standing hypothesis implies that
\(\mu\) is sigma-finite: the closed disks \(\{|w|\leq 1-1/n\}\) have finite
\(\mu\)-measure and exhaust \(\D\).  Hence, for every \(M>0\), there is a
measurable set \(E\subset\D\) with \(0<\mu(E)<\infty\) and \(\mu(E)>M\).
Using \eqref{eq:green-average}, \eqref{eq:boundary-limits}, and the uniform
bound in \eqref{eq:local-limits}, dominated convergence gives
\[
 \lim_{r\uparrow1}\int_E \mathcal D_{\sigma_r(w)}(K)\,d\mu(w)
 =
 \frac{\pi^2}{2}\mu(E).
\]
Therefore the \(M(\mathcal D_\mu)\)-seminorm of \(K=\HH(1)\) is at least
\((\pi/\sqrt2)\sqrt{\mu(E)}\).  Since \(M\) was arbitrary, this seminorm is
infinite.  Thus \(\HH(1)\notin M(\mathcal D_\mu)\), and \(\HH\) is not bounded
from \(H^\infty\) into \(M(\mathcal D_\mu)\).

Now suppose that \(\mu(\D)<\infty\).  Applying
\eqref{eq:green-average}--\eqref{eq:boundary-limits} directly over all of
\(\D\), again by dominated convergence, gives
\begin{equation}\label{eq:L-seminorm}
\sup_{\phi\in\Aut}
\left(\int_{\D}|L'(z)|^2U_\mu(\phi(z))\dA(z)\right)^{1/2}
=\frac{\pi}{\sqrt2}\sqrt{\mu(\D)}.
\end{equation}
The same identity holds with \(L\) replaced by \(K\).

For \(f\in H^\infty\), Bellavita--Stylogiannis prove
\[
 (1-z)(\HH f)'(z)=b_f(z),
 \qquad b_f\in H^\infty,\quad \|b_f\|_\infty\leq\|f\|_\infty .
\]
Also,
\[
 \HH f(0)=\int_0^1 f(t)\,dt.
\]
Using \eqref{eq:L-seminorm},
\[
 \|\HH f\|_{\MD}
 \leq
 \left(1+\frac{\pi}{\sqrt2}\sqrt{\mu(\D)}\right)\|f\|_\infty .
\]
Thus the displayed constant is an upper bound for the operator norm.  Taking
\(f\equiv1\) gives \(\HH f=K\), \(K(0)=1\), and the \(K\)-version of
\eqref{eq:L-seminorm}, so the same constant is attained.  This proves the
formula and the theorem.

\section*{Verification notes and limitations}

The proof is for the \(U_\mu\)-based \(M(\mathcal D_\mu)\) norm displayed in
the source paper.  The most important reviewer checks are:

\begin{itemize}
\item the automorphism-invariance identity for the Green kernel;
\item the boundary-local limits \eqref{eq:boundary-limits}, especially for
  \(K=\HH(1)\);
\item the sigma-finite truncation argument in the infinite-mass case.
\end{itemize}

This packet supersedes the earlier finite-measure partial packet.  The theorem
also makes the finite-mass hypothesis unavoidable for boundedness in this
norm convention.

The local run indexes had only the earlier finite-measure partial packet for
arXiv:2408.01060.  A bounded web search on 2026-06-18 for the arXiv id, the
paper title, ``Question 8.1'', ``non-radial'', ``Hilbert matrix'', and
\(M(\mathcal D_\mu)\) found the source paper and nearby Hilbert-matrix
papers/surveys, but no later exact answer to Question 8.1.

\section*{Human-review recommendation}

Review as a likely valid full answer to Question 8.1 under the source's
displayed \(M(\mathcal D_\mu)\) norm.  The subtle point is compatibility with
the source's later examples/norm conventions for \(Q_p\)-type spaces; the
packet's proof uses the \(U_\mu\)-Green-potential norm exactly as displayed in
the question's section.

\begin{thebibliography}{9}
\bibitem{BellavitaStylogiannis}
C.~Bellavita and G.~Stylogiannis,
\emph{Hilbert matrix operator Acting between conformally invariant spaces},
arXiv:2408.01060, 2024.
\end{thebibliography}

\end{document}
